\documentclass[report, 12pt]{uftreport}

\usepackage{amsmath}

\title{Laboratório 2 -- Display de 7 Segmentos e Conversor Gray-Binário}

\author{Pedro Lucas}{Moreira}

\advisor{Prof.}{Tiago}{Almeida}{}

\department{CC}
\class{Linguagens de Descrição de Hardware -- 2026/2}

\date{8}{9}{2026}

\keyword{VHDL}
\keyword{Display de 7 Segmentos}
\keyword{Decodificador}
\keyword{Atribuição Concorrente}

\begin{document}

\frontmatter
\maketitle
\makefrontpage

\mainmatter
\ChapterStart{first}{Questão 1 -- Display de 7 Segmentos}

\chapter{Questão 1 -- Display de 7 Segmentos}

\section{Item (a) -- Decodificador decimal com detecção de erro}

Cada bit de \texttt{SWITCH\_PINS} corresponde a um switch físico da DE1 (0 = chave para baixo, 1 = chave para cima), formando um número binário de 0 a 127. O decodificador mapeia esse valor diretamente para os segmentos ativos-baixo do display: as combinações de 0 a 9 acendem o dígito correspondente, e qualquer valor fora dessa faixa (10 a 127) acende o padrão da letra E.

\begin{lstlisting}[language=VHDL]
library ieee;
use ieee.std_logic_1164.all;

entity display is
    port (
        SWITCH_PINS: in std_logic_vector(6 downto 0);
        HEX_PINS: out std_logic_vector(6 downto 0)
    );
end entity display;

architecture rtl of display is
begin
    with SWITCH_PINS select
        HEX_PINS <= "1000000" when "0000000", -- 0
                    "1111001" when "0000001", -- 1
                    "0100100" when "0000010", -- 2
                    "0110000" when "0000011", -- 3
                    "0011001" when "0000100", -- 4
                    "0010010" when "0000101", -- 5
                    "0000010" when "0000110", -- 6
                    "1111000" when "0000111", -- 7
                    "0000000" when "0001000", -- 8
                    "0010000" when "0001001", -- 9
                    "0000110" when others;    -- E (fora da faixa 0-9)

end architecture rtl;
\end{lstlisting}

\section{Item (b) -- Conversor hexadecimal}

Para o item (b), a faixa de entrada de 4 bits (0 a 15) é totalmente válida em hexadecimal, então não existe caso de erro: cada combinação binária tem um dígito ou letra (0--9, A--F) correspondente nos segmentos.

\begin{lstlisting}[language=VHDL]
library ieee;
use ieee.std_logic_1164.all;

entity conv_7seg is
    port (
        HEX_IN  : in std_logic_vector(3 downto 0);
        HEX_OUT : out std_logic_vector(6 downto 0)
    );
end entity conv_7seg;

architecture rtl of conv_7seg is
begin
    with HEX_IN select
        HEX_OUT <= "1000000" when "0000", -- 0
                   "1111001" when "0001", -- 1
                   "0100100" when "0010", -- 2
                   "0110000" when "0011", -- 3
                   "0011001" when "0100", -- 4
                   "0010010" when "0101", -- 5
                   "0000010" when "0110", -- 6
                   "1111000" when "0111", -- 7
                   "0000000" when "1000", -- 8
                   "0010000" when "1001", -- 9
                   "0001000" when "1010", -- A
                   "0000011" when "1011", -- b
                   "1000110" when "1100", -- C
                   "0100001" when "1101", -- d
                   "0000110" when "1110", -- E
                   "0001110" when "1111", -- F
                   "1111111" when others;

end architecture rtl;
\end{lstlisting}

O topo do projeto (\texttt{displayb.vhd}) apenas instancia o componente \texttt{conv\_7seg}, ligando as entradas físicas dos switches à sua porta, para teste na DE1:

\begin{lstlisting}[language=VHDL, label={lst:displayb}]
library ieee;
use ieee.std_logic_1164.all;

entity displayb is
    port (
        SWITCH_PINS: in std_logic_vector(3 downto 0);
        HEX_PINS: out std_logic_vector(6 downto 0)
    );
end entity displayb;

architecture rtl of displayb is
begin
    conv_inst: entity work.conv_7seg
        port map (
            HEX_IN  => SWITCH_PINS,
            HEX_OUT => HEX_PINS
        );

end architecture rtl;
\end{lstlisting}

\ChapterStart{second}{Questão 2 -- Conversor Gray -- Binário}

\chapter{Questão 2 -- Conversor Gray -- Binário}

\section{Item (a)}

Como cada bit de saída depende apenas dos bits de entrada (nunca de outro bit de saída), a implementação é uma atribuição concorrente direta, sem necessidade de sinal intermediário nem de processo.

\begin{lstlisting}[language=VHDL]
library ieee;
use ieee.std_logic_1164.all;

entity gray2bin is
    port (
        GRAY_IN : in  std_logic_vector(3 downto 0);
        BIN_OUT : out std_logic_vector(3 downto 0)
    );
end entity gray2bin;

architecture rtl of gray2bin is
begin
    BIN_OUT(3) <= GRAY_IN(3);
    BIN_OUT(2) <= GRAY_IN(2) xor GRAY_IN(3);
    BIN_OUT(1) <= GRAY_IN(1) xor GRAY_IN(2) xor GRAY_IN(3);
    BIN_OUT(0) <= GRAY_IN(0) xor GRAY_IN(1) xor GRAY_IN(2) xor GRAY_IN(3);

end architecture rtl;
\end{lstlisting}

\section{Itens (d) e (e)}

O \texttt{demo\_setup} foi feito em VHDL, instanciando os dois componentes já prontos \texttt{gray2bin} e \texttt{conv\_7seg} e ligando-os em cascata: a entrada de 4 bits em código Gray (\texttt{SW}) é convertida para binário e em seguida decodificada para os 7 segmentos (\texttt{HEX3}).

\begin{lstlisting}[language=VHDL]
library ieee;
use ieee.std_logic_1164.all;

entity demo_setup is
    port (
        SW    : in  std_logic_vector(3 downto 0);
        HEX3  : out std_logic_vector(6 downto 0)
    );
end entity demo_setup;

architecture rtl of demo_setup is
    signal bin_value : std_logic_vector(3 downto 0);
begin
    gray_inst: entity work.gray2bin
        port map (
            GRAY_IN => SW,
            BIN_OUT => bin_value
        );

    seg_inst: entity work.conv_7seg
        port map (
            HEX_IN  => bin_value,
            HEX_OUT => HEX3
        );

end architecture rtl;
\end{lstlisting}

\section{Item (f)}

A atribuição de pinos foi feita via script Tcl, cobrindo as 4 entradas de \texttt{SW} e os 7 segmentos de \texttt{HEX3} conforme a tabela oficial de pinos da DE1:

\begin{lstlisting}[language=tcl]
package require ::quartus::project
package require ::quartus::flow

set project_name "demo_setup"

project_new $project_name -overwrite

set_global_assignment -name FAMILY "Cyclone II"
set_global_assignment -name DEVICE EP2C20F484C7
set_global_assignment -name VHDL_FILE a.vhd
set_global_assignment -name VHDL_FILE ../01/conv_7seg.vhd
set_global_assignment -name VHDL_FILE demo_setup.vhd
set_global_assignment -name TOP_LEVEL_ENTITY demo_setup

set_location_assignment PIN_L22 -to SW[0]
set_location_assignment PIN_L21 -to SW[1]
set_location_assignment PIN_M22 -to SW[2]
set_location_assignment PIN_V12 -to SW[3]

set_location_assignment PIN_F4 -to HEX3[0]
set_location_assignment PIN_D5 -to HEX3[1]
set_location_assignment PIN_D6 -to HEX3[2]
set_location_assignment PIN_J4 -to HEX3[3]
set_location_assignment PIN_L8 -to HEX3[4]
set_location_assignment PIN_F3 -to HEX3[5]
set_location_assignment PIN_D4 -to HEX3[6]

export_assignments
execute_flow -compile

project_close
\end{lstlisting}

\end{document}
